\scanpage{138}
% archive foliation: 137
\[
z_1=u(e_1),\ z_2=u(e_2)\in H\subset V
\]
alors la structure de Torelli $\uline{u\circ g}$ correspond à $z'_1,z'_2\in\Pi$ donnés par
\[
\begin{aligned}
z'_1&=ug(e_1)=u(ae_1+ce_2)=az_1+cz_2\\
z'_2&=ug(e_2)=u(be_1+de_2)=bz_1+dz_2,
\end{aligned}
\]
qui correspond donc à un élément
\[
z'=x'+iy'\in\mathcal D'
\]
donné par
\deleted{\unsure{(24)}}
\begin{equation}
z'\deleted{\unsure{(x'+iy')}}=
-\frac{z'_2}{z'_1}
=-\frac{bz_1+dz_2}{az_1+cz_2}
=-\frac{d\frac{z_2}{z_1}+b}{c\frac{z_2}{z_1}+a}
=\frac{d\unsure{z|\overline z}\unsure{-|+}b}{-c\unsure{z|\overline z}+a}
\tag{25}
\end{equation}
\deleted{$\unsure{\widetilde z'=-\overline{z'}}$}

où $\deleted{\unsure{\widetilde z=x+iy=-\overline z}}$.
\margin{\diagram{\begin{tikzcd}[ampersand replacement=\&,column sep=small,row sep=large]
 \& z'=\unsure{\frac{-dz+b}{-cz+a}} \arrow[dl,bend right=55] \arrow[d,bend left=50] \arrow[d,bend right=25] \\
\unsure{(?)} \& z'=\unsure{\frac{dz-b}{-cz+a}} \arrow[l,bend left=35] \arrow[u,bend left=35]
\end{tikzcd}}}

Si on \deleted{\unsure{veut}} préfère expliciter l'action à gauche $u\longmapsto u\circ g^{-1}$ de $\mathrm{Sl}(2,\mathbb Z)$ sur l'ens. des structures de Torelli mod iso, on voit que
\[
g^{-1}=\begin{pmatrix}d&-b\\-c&a\end{pmatrix}
\qquad\deleted{\text{d'où}}
\]
d'où maintenant, pour $z'=x'+iy'$ correspondant à $u\circ g^{-1}$
\deleted{\unsure{(25 bis)}}
\[
\deleted{\widetilde z'=\overline{z'}
=\frac{a\widetilde z-b}{-c\widetilde z+d}
=\frac{-a\overline z-b}{-c\overline z-d}
=\unsure{\frac{a\overline z+b}{c\overline z+d}|\overline{\frac{az+b}{cz+d}}}}
\]
\begin{equation}
z'=\frac{az+b}{cz+d}
\tag{25 bis}
\end{equation}
\uline{victoire}, c'est la formule habituelle,

mais on fera attention que le nb complexe $z\in\mathcal D$ attaché à la structure de Torelli définie par \deleted{$z_1,z_2\in t_E$}\inserted{les deux} éléments $z_1,z_2$ de la droite complexe $\deleted{\unsure{t=\lambda}}\ t_E$ est donné, non
